% =====================================================================
% Muc bo sung cho ReportLatexType1/bao_cao.tex (va ban IEEE).
% Vi tri de xuat: cuoi \section{Ket qua thuc nghiem}, ngay truoc
% \section{Danh gia tinh kha thi cua de xuat}. Doan cuoi (\paragraph{Thoi diem cua du lieu})
% nen chen vao \subsection{Moi de doa doi voi tinh hop le}.
% Ma va ket qua: thu muc cai_tien_S/ (cai_tien_S.py, kiem_tra_ro_ri.py, kiem_tra_cpe.py,
% msrc_kiem_chung.py, cvss_tai_ngay_D.py), chay ngay 14/10/2026.
% =====================================================================

\subsection{Các cải tiến đã thử nhưng không hiệu quả}\label{sec:khonghieuqua}

Giới hạn lớn nhất của $S$ là gần như không cảnh báo trước được lỗ hổng mới công bố
(Bảng~\ref{tab:tuoi}). Nhóm thử ba cách bổ sung tín hiệu có sẵn sớm hơn EPSS. Cả ba dùng đúng
thiết kế của thực nghiệm 4: mỗi tháng công bố là một đợt, ngày quyết định $D$ là cuối tháng, nhãn
là vào KEV trong 180 ngày sau $D$, kiểm thử 2024 và 2025 (84.821 lỗ hổng, 116 bị khai thác). Vì
$D$ là cuối tháng công bố, mọi lỗ hổng trong tập mới công bố tối đa 31 ngày. Với công thức gốc,
thiết kế này cho lại đúng số của Bảng~\ref{tab:a1}. Tham số của các biến thể chỉ học trên những
đợt có cửa sổ nhãn kết thúc trước ngày 01/01 của năm kiểm thử, và mọi tín hiệu chỉ tính nếu có
trước hoặc đúng ngày $D$.

Các biến thể giữ nguyên dạng của công thức, chỉ nâng sàn cho EPSS ở nhánh chưa có trong KEV:
\begin{equation}
S' = \text{CVSS} \times \max\bigl(\text{EPSS},\ p\bigr)^{0{,}3},
\end{equation}
trong đó $p$ là tỷ lệ vào KEV học được của nhóm lỗ hổng có cùng tín hiệu. Lỗ hổng có EPSS cao
giữ nguyên điểm; chỉ lỗ hổng có EPSS gần 0 nhưng đã có tín hiệu khác được nâng lên.

\begin{table}
\centering
\caption{Số lỗ hổng bị khai thác bắt được khi vá cùng số lượng, gộp 2024 và 2025 (116 lỗ hổng
bị khai thác). Chênh lệch so với công thức gốc kèm khoảng tin cậy 95\% bootstrap 1.000 lần.}
\label{tab:khonghieuqua}
\begin{tabular}{lccccc}
\toprule
Cách xếp & 1\% & 5\% & 10\% & 25\% & Bằng CSA \\
\midrule
$S$ gốc & 20,0 & 40,0 & 48,7 & 72,0 & 72,0 \\
$S$ + PoC (Exploit-DB) & 20,0 & 40,0 & 48,7 & 72,0 & 72,0 \\
$S$ + nhà cung cấp theo CPE & 21,9 & 43,6 & 59,1 & 87,4 & 87,5 \\
$S$ + nhà cung cấp theo CNA & 21,4 & 40,4 & 49,1 & 75,8 & 77,0 \\
$S$ + cờ Exploited của Microsoft & 22,0 & 42,0 & 50,6 & 73,0 & 73,0 \\
\midrule
\multicolumn{6}{l}{Chênh lệch so với $S$ gốc} \\
\quad Nhà cung cấp theo CPE & $+1{,}9$ & $+3{,}6$ & $+10{,}4$ & $+15{,}4^{*}$ & $+15{,}5^{*}$ \\
\quad Nhà cung cấp theo CNA & $+1{,}4$ & $+0{,}4$ & $+0{,}4$ & $+3{,}8$ & $+5{,}0$ \\
\quad Cờ Exploited của Microsoft & $+2{,}0$ & $+2{,}0$ & $+1{,}9$ & $+1{,}0$ & $+1{,}0$ \\
\bottomrule
\multicolumn{6}{l}{\footnotesize $^{*}$ khoảng tin cậy không chứa 0; mọi chênh lệch khác đều chứa 0.}
\end{tabular}
\end{table}

\paragraph{PoC công khai.} Thực nghiệm 5 cho thấy điều kiện có PoC giúp quy tắc CSA ở tập lỗ hổng
tồn đọng. Với lỗ hổng mới thì không: tại ngày $D$, chỉ 3 trong 116 lỗ hổng bị khai thác (2,6\%)
đã có PoC trên Exploit-DB. Tỷ lệ vào KEV của nhóm có PoC trên dữ liệu huấn luyện bằng 0 (kiểm thử
2024) và 0,71\% (kiểm thử 2025), quá thấp để thay đổi thứ hạng ở mọi khối lượng vá được xét.
Phần lớn PoC trên Exploit-DB chưa xuất hiện vào cuối tháng công bố: chỉ 0,10\% lỗ hổng trong tập
đã có PoC tại ngày $D$.

\paragraph{Lịch sử của nhà cung cấp.} Lỗ hổng mới chưa có lịch sử, nhưng nhà cung cấp của nó thì
có. Biến thể chia lỗ hổng theo số mục KEV của nhà cung cấp tính đến ngày $D$ (0, 1--4, 5--19,
20--99, từ 100). Khi lấy nhà cung cấp từ CPE trong cấu hình của NVD, biến thể bắt thêm 15,5 ca ở
khối lượng CSA, khoảng tin cậy $[+4{,}0;\ +27{,}9]$. Nhưng phần lớn mức tăng này là rò rỉ: CPE do
NVD gán khi phân tích, thường sau ngày công bố, và 100\% lỗ hổng bị khai thác có CPE trong khi
nhóm còn lại chỉ 81,1\% (2024) và 66,4\% (2025). Lỗ hổng thiếu CPE tự động rơi vào bậc thấp nhất,
và đó gần như toàn là lỗ hổng không bị khai thác. Lấy nhà cung cấp từ CNA (\texttt{sourceIdentifier},
có ngay khi công bố), mức tăng chỉ còn $+5{,}0$ ca, khoảng tin cậy $[-3{,}3;\ +13{,}8]$. Một phần vì
CNA lớn nhất là \texttt{cve@mitre.org}, gộp nhiều nhà cung cấp làm một. Khi chỉ so trên lỗ hổng
đã có CPE (bỏ lợi thế do thiếu CPE), mức tăng là $+13{,}2$ ca $[+1{,}6;\ +25{,}5]$ ở khối lượng
CSA nhưng chỉ $+0{,}9$ ca $[-1{,}9;\ +3{,}2]$ khi vá 1\%. Vậy lịch sử nhà cung cấp có tín hiệu thật
khi vá nhiều, nhưng không giúp nhóm vá khẩn cấp, đúng nơi cần cảnh báo sớm.

\paragraph{Nguồn cảnh báo sớm hơn CISA KEV.} Microsoft đánh dấu \texttt{Exploited:Yes} trong
tài liệu CVRF của MSRC, có ghi ngày trong lịch sử sửa đổi. Trong 108 lỗ hổng có cả hai nguồn,
Microsoft báo cùng ngày với CISA ở 85 trường hợp (79\%), báo trước ở 20 (10 trường hợp chỉ trước
1 ngày) và báo sau ở 3; trung vị độ trễ bằng 0 ngày. Mở rộng thành phần KEV thành
``CISA KEV hoặc Microsoft đã báo bị khai thác'' chỉ bắt thêm 1--2 ca, không có ý nghĩa thống kê.
Với nhà cung cấp lớn nhất, CISA KEV đã gần như theo kịp nguồn gốc; lợi ích của nguồn cảnh báo
sớm, nếu có, nằm ở các nhà cung cấp mà CISA cập nhật chậm.

\paragraph{Nhận xét.} Không cách nào bù được giới hạn với lỗ hổng mới công bố. Tín hiệu có ích
nhất (lịch sử nhà cung cấp) chỉ cải thiện thứ tự khi vá nhiều, và một tổ chức có thể dùng nó trong
chu kỳ vá thông thường nhờ danh mục tài sản của chính mình. Với nhóm vá khẩn cấp, kết luận ở
Mục~\ref{sec:khathi} giữ nguyên: phải dựa vào quy tắc KEV, biện pháp giảm thiểu tạm thời cho hệ
thống mở ra Internet và tốc độ phản ứng, không dựa vào dự báo.

% ---------------------------------------------------------------------
% Chen vao \subsection{Moi de doa doi voi tinh hop le}
% ---------------------------------------------------------------------
\paragraph{Thời điểm của dữ liệu.} Các thực nghiệm lấy điểm CVSS do NVD chấm khi có, nếu không
thì lấy điểm khác (của CNA hoặc CISA-ADP), và bỏ lỗ hổng không có điểm. Điểm NVD và điểm của
CISA-ADP thường được bổ sung sau ngày công bố, nên chưa chắc có tại ngày quyết định. NVD chấm cao
hơn CNA ở 49,5\% số lỗ hổng có cả hai điểm và thấp hơn ở 27,8\%. Năm 2025, 100\% lỗ hổng KEV có
điểm, nhóm còn lại chỉ 94,2\%. Bảng~\ref{tab:cvssngayD} chạy lại thực nghiệm 4 với bốn cách chọn
điểm. Khi chỉ dùng điểm CNA (có ngay khi công bố), lợi thế của $S$ so với CVSS khi vá 5\% giảm từ
$+18{,}2$ ca $[+7{,}5;\ +28{,}8]$ xuống $+5{,}3$ ca $[-6{,}2;\ +16{,}0]$; còn quy tắc CSA chỉ chọn
11,7\% lỗ hổng và bỏ sót 57,8\% lỗ hổng bị khai thác. Như vậy kết luận ``quy tắc CSA bỏ sót nhiều''
còn mạnh hơn khi dùng điểm có tại thời điểm quyết định, nhưng một phần lợi thế của $S$ so với CVSS
có thể đến từ điểm NVD chấm muộn.

\begin{table}
\centering
\caption{Thực nghiệm 4 với bốn cách chọn điểm CVSS, gộp 2024 và 2025. Cột cuối là số ca $S$ bắt
được nhiều hơn CVSS khi vá 5\%, kèm khoảng tin cậy 95\%.}\label{tab:cvssngayD}
\begin{tabular}{lrrrc}
\toprule
Cách chọn điểm & Lỗ hổng & CSA phải vá & CSA bỏ sót & $S$ hơn CVSS, vá 5\% \\
\midrule
NVD trước (như các thực nghiệm) & 84.821 & 25,1\% & 43,1\% & $+18{,}2$ $[+7{,}5;\ +28{,}8]$ \\
Nguồn sớm trước (CNA, ADP, NVD) & 84.821 & 20,2\% & 37,9\% & $+12{,}4$ $[+1{,}2;\ +24{,}2]$ \\
Chỉ CNA, giữ lỗ hổng không có điểm & 87.910 & 11,7\% & 57,8\% & $+5{,}3$ $[-6{,}2;\ +16{,}0]$ \\
Chỉ CNA, bỏ lỗ hổng không có điểm & 56.985 & 18,0\% & 33,8\% & $+1{,}1$ $[-8{,}8;\ +10{,}9]$ \\
\bottomrule
\end{tabular}
\end{table}
